\section*{ანოტაცია}
\addcontentsline{toc}{section}{ანოტაცია}

ერთობლივი მუშაობა ერთსა და იმავე ტექსტურ ფაილზე დღეს ყოველდღიურობის ნაწილია
- იქნება ეს წყვილში პროგრამირება, სემინარი, ლაბორატორიული სამუშაო თუ
ტექნიკური გასაუბრება. ამის მიუხედავად, არსებული ინსტრუმენტების დიდი ნაწილი
ორ დაშვებას ეყრდნობა: დოკუმენტი უნდა „აიტვირთოს“ რომელიღაც კომპანიის
ცენტრალურ სერვერზე (მაგ. Google Docs) და ყველა მონაწილეს სტაბილური ინტერნეტი უნდა ჰქონდეს.
ჩვენი გამოცდილებით, უნივერსიტეტის აუდიტორიაში ორივე დაშვება ხშირად
ირღვევა: ეკრანის გაზიარება ხშირად არ მუშაობს და სემინარზე
ყველა სტუდენტს პროფესორის ლეპტოპის გარშემო შეკრება უწევს.

ამ პრობლემის გადასაჭრელად შევქმენით Peared, კროსპლატფორმული desktop
აპლიკაცია ერთობლივი ტექსტური რედაქტირებისთვის, რომელიც მუშაობს როგორც
ლოკალურად, ისე გლობალურად. სესიის
წამომწყები (ჰოსტი) საკუთარ კომპიუტერზე უშვებს მსუბუქ relay სერვერს,
რომელსაც აპლიკაცია ავტომატურად მართავს, ხოლო დანარჩენი მონაწილეები
უბრალოდ ბმულით უერთდებიან. დოკუმენტი არასდროს ტოვებს მონაწილეთა
კომპიუტერებს. 

პარალელური რედაქტირებისას ცვლილებების უკონფლიქტოდ შერწყმას უზრუნველყოფს
CRDT (Conflict-free Replicated Data Type) ალგორითმი, კონკრეტულად, YATA
ვარიანტი, რომელიც კოლაბორაციული ედიტინგის ყველაზე ცნობილი ბიბლიოთეკის,
Yjs-ის, საფუძველია. თითოეული კლიენტი დოკუმენტის
საკუთარ ასლს ფლობს და ოპერაციები ისეა აგებული, რომ ნებისმიერი თანმიმდევრობით
მიღების შემთხვევაში ყველა მონაწილე ერთსა და იმავე ტექსტს დაინახავს. სისტემა
სამი დამოუკიდებელი კომპონენტისგან შედგება: Rust-ზე დაწერილი CRDT ბიბლიოთეკა,
Tauri Framework-ზე აგებული desktop კლიენტი (React/Monaco ფრონტენდით) და Go-ზე
დაწერილი WebSocket relay სერვერი. გლობალური წვდომისთვის აპლიკაცია
ავტომატურად ხსნის Cloudflare Tunnel-ს, რაც პორტების გახსნისა და ქსელის
კონფიგურაციის გარეშე აძლევს დაშორებულ მონაწილეებს შეერთების საშუალებას.

შედეგად მივიღეთ მომუშავე პროდუქტი, რომლითაც რამდენიმე ადამიანს
ერთდროულად შეუძლია ერთი და იმავე ფაილის რედაქტირება და უფლებების მართვა,
და რომელიც ნელი ინტერნეტის პირობებშიც კი ეფექტურად მუშაობს.

\clearpage
\section*{Abstract}
\addcontentsline{toc}{section}{Abstract}

Working together on the same text file is part of everyday life today -
whether it is pair programming, a seminar, a lab session or a technical
interview. Despite this, most existing tools rest on two assumptions: the
document must be uploaded to some company's central server (e.g. Google
Docs), and every participant must have a stable internet connection. In
our experience, both assumptions frequently break down in a university
classroom: screen sharing often fails, and the whole seminar ends up
crowding around the professor's laptop.

To solve this problem we built Peared, a cross-platform desktop
application for collaborative text editing that works both locally and
globally. The session initiator (the host) runs a lightweight relay
server on their own machine, managed automatically by the application,
while the other participants simply join via a link. The document never
leaves the participants' computers.

Conflict-free merging of concurrent edits is guaranteed by a CRDT
(Conflict-free Replicated Data Type) algorithm, specifically the YATA
variant that underlies Yjs, the best-known collaborative editing library.
Every client owns its own replica of the document, and operations are
constructed so that, regardless of the order in which they are received,
all participants will see the same text. The system consists of three
independent components: a CRDT library written in Rust, a desktop client
built on the Tauri framework (with a React/Monaco frontend), and a
WebSocket relay server written in Go. For global access the application
automatically opens a Cloudflare Tunnel, which lets remote participants
connect without port forwarding or any network configuration.

The result is a working product with which several people can edit the
same file simultaneously and manage write permissions, and which performs
well even under a slow internet connection.
